% GENERATED FILE. Edit canonical lessons, then run npm run generate:books.
% canonical-source-hash: fnv1a64:da778ae131956da6
% canonical-lessons: JA-C454-torikumu, JA-C454-hakadoru, JA-C454-hikiukeru, JA-C454-yarinaosu, JA-C454-shippai-suru

\chapter{Getting the Job Done}
\label{ch:ja-getting-the-job-done}
\hlchaptermodality{454}

\begin{chapteropening}
\textbf{By the end of this chapter:} \emph{I can say how a job is going.}

Complete the last lesson of chapter 454: I can say how a job is going.
\end{chapteropening}

\section[\texorpdfstring{\ja{とりくむ} (torikumu)}{とりくむ (torikumu)}]{\ja{とりくむ} (torikumu) — to tackle; to work hard on}

\label{lesson:JA-C454-torikumu}



\begin{quote}
Before the new one: say the Japanese for to hand in, then the Japanese for to stay up all night.
\end{quote}

\subsection*{You'll want to know: \ja{とりくむ}}
\textbf{\ja{とりくむ}} — \emph{torikumu} — \textquotedblleft{}to tackle; to work hard on\textquotedblright{}.

\begin{quote}
\textbf{\ja{みんなで} \ja{もんだいに} \ja{とりくみました。}} — \emph{minna de mondai ni torikumimashita} — \textquotedblleft{}We all tackled the problem together.\textquotedblright{}
\end{quote}

Polite \textbf{\ja{とりくみます}}, \emph{torikumimasu}; a godan verb. Past (polite): \textbf{\ja{せんせいが} \ja{もんだいに} \ja{とりくみました。}}, \emph{sensei ga mondai ni torikumimashita}.

\begin{cousinweb}
\textbf{\ja{とりくむ}} (\emph{torikumu}) is a native Japanese compound (wago), inherited rather than borrowed from Chinese: \emph{tori}, \textquotedblleft{}taking (toru)\textquotedblright{} and \emph{kumu}, \textquotedblleft{}putting together\textquotedblright{} (JMdict, KANJIDIC2).
\end{cousinweb}

\begin{grammarlens}[title={What you've built}]
One more word for describing what you lived and felt.
\end{grammarlens}

\subsection*{Your turn}
\begin{itemize}
\raggedright
  \item \emph{Say it:} \emph{torikumu}
  \item \emph{Say it:} \emph{torikumu}, once more
  \item \emph{Say it:} \emph{minna de mondai ni torikumimashita} — We all tackled the problem together.
  \item \emph{Recall:} say the Japanese for to hand in, then the Japanese for to stay up all night, then say \emph{torikumu} again
\end{itemize}

\begin{tcolorbox}[breakable,colback=teal!4,colframe=teal!35!black,title={Before you move on}]
What does \ja{とりくむ} mean? (\textquotedblleft{}To tackle; to work hard on\textquotedblright{}.) Say it once more.
\end{tcolorbox}
\section[\texorpdfstring{\ja{はかどる} (hakadoru)}{はかどる (hakadoru)}]{\ja{はかどる} (hakadoru) — to make good progress}

\label{lesson:JA-C454-hakadoru}



\begin{quote}
Before the new one: say the Japanese for to stay up all night, then the Japanese for to tackle.
\end{quote}

\subsection*{You'll want to know: \ja{はかどる}}
\textbf{\ja{はかどる}} — \emph{hakadoru} — \textquotedblleft{}to make good progress\textquotedblright{}.

\begin{quote}
\textbf{\ja{けさは} \ja{しごとが} \ja{はかどりました。}} — \emph{kesa wa shigoto ga hakadorimashita} — \textquotedblleft{}My work went well this morning.\textquotedblright{}
\end{quote}

Polite \textbf{\ja{はかどります}}, \emph{hakadorimasu}; a godan verb. Past (polite): \textbf{\ja{べんきょうが} \ja{はかどりました。}}, \emph{benkyō ga hakadorimashita}.

\begin{cousinweb}
\textbf{\ja{はかどる}} (\emph{hakadoru}) is a native Japanese word (wago), inherited rather than borrowed from Chinese; its kanji spelling was attached to it later (JMdict).
\end{cousinweb}

\begin{grammarlens}[title={What you've built}]
One more word for describing what you lived and felt.
\end{grammarlens}

\subsection*{Your turn}
\begin{itemize}
\raggedright
  \item \emph{Say it:} \emph{hakadoru}
  \item \emph{Say it:} \emph{hakadoru}, once more
  \item \emph{Say it:} \emph{kesa wa shigoto ga hakadorimashita} — My work went well this morning.
  \item \emph{Recall:} say the Japanese for to stay up all night, then the Japanese for to tackle, then say \emph{hakadoru} again
\end{itemize}

\begin{tcolorbox}[breakable,colback=teal!4,colframe=teal!35!black,title={Before you move on}]
What does \ja{はかどる} mean? (\textquotedblleft{}To make good progress\textquotedblright{}.) Say it once more.
\end{tcolorbox}
\section[\texorpdfstring{\ja{ひきうける} (hikiukeru)}{ひきうける (hikiukeru)}]{\ja{ひきうける} (hikiukeru) — to take on (a task)}

\label{lesson:JA-C454-hikiukeru}



\begin{quote}
Before the new one: say the Japanese for to tackle, then the Japanese for to make good progress.
\end{quote}

\subsection*{You'll want to know: \ja{ひきうける}}
\textbf{\ja{ひきうける}} — \emph{hikiukeru} — \textquotedblleft{}to take on (a task)\textquotedblright{}.

\begin{quote}
\textbf{\ja{かれは} \ja{ぶちょうの} \ja{しごとを} \ja{ひきうけました。}} — \emph{kare wa buchō no shigoto o hikiukemashita} — \textquotedblleft{}He took on the department head's work.\textquotedblright{}
\end{quote}

Polite \textbf{\ja{ひきうけます}}, \emph{hikiukemasu}; an ichidan verb. Past (polite): \textbf{\ja{どうりょうが} \ja{しごとを} \ja{ひきうけました。}}, \emph{dōryō ga shigoto o hikiukemashita}.

\begin{cousinweb}
\textbf{\ja{ひきうける}} (\emph{hikiukeru}) is a native Japanese compound (wago), inherited rather than borrowed from Chinese: \emph{hiki}, \textquotedblleft{}pulling (hiku)\textquotedblright{} and \emph{ukeru}, \textquotedblleft{}to receive\textquotedblright{} (JMdict, KANJIDIC2).
\end{cousinweb}

\begin{grammarlens}[title={What you've built}]
One more word for describing what you lived and felt.
\end{grammarlens}

\subsection*{Your turn}
\begin{itemize}
\raggedright
  \item \emph{Say it:} \emph{hikiukeru}
  \item \emph{Say it:} \emph{hikiukeru}, once more
  \item \emph{Say it:} \emph{kare wa buchō no shigoto o hikiukemashita} — He took on the department head's work.
  \item \emph{Recall:} say the Japanese for to tackle, then the Japanese for to make good progress, then say \emph{hikiukeru} again
\end{itemize}

\begin{tcolorbox}[breakable,colback=teal!4,colframe=teal!35!black,title={Before you move on}]
What does \ja{ひきうける} mean? (\textquotedblleft{}To take on (a task)\textquotedblright{}.) Say it once more.
\end{tcolorbox}
\section[\texorpdfstring{\ja{やりなおす} (yarinaosu)}{やりなおす (yarinaosu)}]{\ja{やりなおす} (yarinaosu) — to do over; to redo}

\label{lesson:JA-C454-yarinaosu}



\begin{quote}
Before the new one: say the Japanese for to make good progress, then the Japanese for to take on.
\end{quote}

\subsection*{You'll want to know: \ja{やりなおす}}
\textbf{\ja{やりなおす}} — \emph{yarinaosu} — \textquotedblleft{}to do over; to redo\textquotedblright{}.

\begin{quote}
\textbf{\ja{まちがえたので、さいしょから} \ja{やりなおしました。}} — \emph{machigaeta node, saisho kara yarinaoshimashita} — \textquotedblleft{}I made a mistake, so I did it over from the start.\textquotedblright{}
\end{quote}

Polite \textbf{\ja{やりなおします}}, \emph{yarinaoshimasu}; a godan verb. Past (polite): \textbf{\ja{がくせいが} \ja{しゅくだいを} \ja{やりなおしました。}}, \emph{gakusei ga shukudai o yarinaoshimashita}.

\begin{cousinweb}
\textbf{\ja{やりなおす}} (\emph{yarinaosu}) is a native Japanese compound (wago), inherited rather than borrowed from Chinese: \emph{yari}, \textquotedblleft{}doing (yaru)\textquotedblright{} and \emph{naosu}, \textquotedblleft{}to fix\textquotedblright{} (JMdict, KANJIDIC2).
\end{cousinweb}

\begin{grammarlens}[title={What you've built}]
One more word for describing what you lived and felt.
\end{grammarlens}

\subsection*{Your turn}
\begin{itemize}
\raggedright
  \item \emph{Say it:} \emph{yarinaosu}
  \item \emph{Say it:} \emph{yarinaosu}, once more
  \item \emph{Say it:} \emph{machigaeta node, saisho kara yarinaoshimashita} — I made a mistake, so I did it over from the start.
  \item \emph{Recall:} say the Japanese for to make good progress, then the Japanese for to take on, then say \emph{yarinaosu} again
\end{itemize}

\begin{tcolorbox}[breakable,colback=teal!4,colframe=teal!35!black,title={Before you move on}]
What does \ja{やりなおす} mean? (\textquotedblleft{}To do over; to redo\textquotedblright{}.) Say it once more.
\end{tcolorbox}
\section[\texorpdfstring{\ja{しっぱいする} (shippai suru)}{しっぱいする (shippai suru)}]{\ja{しっぱいする} (shippai suru) — to fail}

\label{lesson:JA-C454-shippai-suru}



\begin{quote}
Before the new one: say the Japanese for to take on, then the Japanese for to do over.
\end{quote}

\subsection*{You'll want to know: \ja{しっぱいする}}
\textbf{\ja{しっぱいする}} — \emph{shippai suru} — \textquotedblleft{}to fail\textquotedblright{}.

\begin{quote}
\textbf{\ja{かれは} \ja{しけんに} \ja{しっぱいしました。}} — \emph{kare wa shiken ni shippai shimashita} — \textquotedblleft{}He failed the exam.\textquotedblright{}
\end{quote}

Polite \textbf{\ja{しっぱいします}}, \emph{shippaishimasu}; a \ja{する} verb. Past (polite): \textbf{\ja{せんしゅが} \ja{しっぱいしました。}}, \emph{senshu ga shippai shimashita}.

\begin{cousinweb}
\textbf{\ja{しっぱいする}} (\emph{shippai suru}) is Sino-Japanese (kango) \emph{shippai} with the native verb \emph{suru}, \textquotedblleft{}to do\textquotedblright{}: two Chinese characters read with their on'yomi, readings borrowed from Chinese. They are \emph{shi}, \textquotedblleft{}lose\textquotedblright{} (Middle Chinese \emph{syit}) and \emph{pai}, \textquotedblleft{}failure\textquotedblright{} (\emph{baejH}) (KANJIDIC2; Middle Chinese in Baxter's transcription).
\end{cousinweb}

\begin{grammarlens}[title={What you've built}]
One more word for describing what you lived and felt.
\end{grammarlens}

\subsection*{Your turn}
\begin{itemize}
\raggedright
  \item \emph{Say it:} \emph{shippai suru}
  \item \emph{Say it:} \emph{shippai suru}, once more
  \item \emph{Say it:} \emph{kare wa shiken ni shippai shimashita} — He failed the exam.
  \item \emph{Recall:} say the Japanese for to take on, then the Japanese for to do over, then say \emph{shippai suru} again
\end{itemize}

\begin{tcolorbox}[breakable,colback=teal!4,colframe=teal!35!black,title={Before you move on}]
What does \ja{しっぱいする} mean? (\textquotedblleft{}To fail\textquotedblright{}.) Say it once more.
\end{tcolorbox}
